The Prover-Verifier game framework has been shown to be useful in training legible language models. However, could we translate the same formalization for AI Control? In this section, we develop a control-robust variant that preserves the theoretical guarantees of the original while addressing its vulnerabilities to strategic deception.

\section{Preliminaries}
We work with measurable spaces $(\mathcal{X}, \mathcal{B}_X)$ and $(\mathcal{Z}, \mathcal{B}_Z)$ representing problems and solutions respectively, with a fixed distribution $P_X$ over problems. For every $x\in\mathcal{X}$ we assume \emph{instance solvability}

\[
  \forall x\;\exists z\in\mathcal{Z}:\;C(x,z)=1,
\]

where $C:\mathcal{X}\times\mathcal{Z}\to\{0,1\}$ is an externally
supplied checker.
\\
The game involves two entities: a \emph{prover} and a \emph{verifier}. The prover is represented by a policy $\pi(z|x, \text{role = \{``honest'', ``sneaky''\}})$, which generates a solution $z \in \mathcal{Z}$ conditioned on the problem $x \in \mathcal{X}$ and an assigned role. The role determines the prover's immediate objective. Specifically, the role could be:
\begin{itemize}
    \item \textbf{Honest.} The prover, denoted $\pi_H(x) = \pi(x, \text{role=``honest''})$, aims to produce a correct solution that the verifier considers convincing.
    \item \textbf{Sneaky.} The prover, denoted $\pi_S(x) = \pi(x, \text{role=``sneaky''})$, aims to produce an incorrect solution that nevertheless convinces the verifier\footnote{For our purpose, it is irrelevant whether the prover policy is unique and role conditioned --namely, the parameter space is shared among the two modes-- or whether the two roles make use of separate copies of the same initial policy. While shared policies are common in symmetric games \citep{silver2017alphazero}, if the game itself is fundamentally asymmetric (e.g., one player is a ``defender'' with different rules/objectives than an ``attacker'', or different starting conditions that aren't just a color swap), then separate networks might be essential because the optimal strategies are inherently different.}.
\end{itemize}
Both $\pi_H(x)$ and $\pi_S(x)$ are assumed to be drawn from the same underlying model class, denoted $\mathcal{P}$. Theoretically, $\mathcal{P}$ is considered to be the set of all deterministic mappings from $\mathcal{X}$ to $\mathcal{Z}$, although in practice we model $\mathcal{H}$ as a Large Language Model (LLM). Although the proofs assume deterministic policies, an LLM with sampling can be viewed as a randomized policy $\tilde{\pi}(x, u) \mapsto z$, where $u$ is a private seed. Modelling it as an input, or fixing it, and because every theorem is stated in expectation over the probability space $(\Omega,\mathcal F,\mathbb P)$, we may absorb $u$ into $\Omega$ and treat $\tilde{\pi}$ as deterministic in its enlarged argument list. Thus all deterministic-policy results apply to stochastic LLMs without modification.\todo[Joan]{This is quite a simplification. Doublecheck with Alessandro/Charlie if it makes sense.}
\\
The verifier $V:\mathcal{X}\times\mathcal{Z}\to[-B,B]$ outputs a
(real-valued) \emph{convincingness score} and is similarly modelled as an LLM.  All random variables live on one probability space $(\Omega,\mathcal F,\mathbb P)$.
\\
A fundamental requirement for this game, following \citet{kirchner2024prover}, is \emph{checkability}—the verifier's capacity to independently assess the correctness of prover-generated solutions. Checkability comprises two essential properties:

\begin{enumerate}
    \item \textbf{Completeness:} The verifier correctly accepts all valid solutions from the honest prover. Formally:
 $$\corr(x, \pi(x)) = v(x, \pi(x)) = 1, \qquad \forall x \in \mathcal{X}$$

    \item \textbf{Soundness:} The verifier correctly rejects all invalid solutions, regardless of which prover in the model class generated them. Formally:
$$\corr(x, \pi'(x)) = 0 \implies v(x, \pi'(x)) = 0, \qquad \forall \pi' \in \mathcal{H}, \forall x \in \mathcal{X}$$
\end{enumerate}
\noindent~These definitions draw inspiration from the literature on Probably-Approximately Correct (PAC) verification and interactive proof systems \citep{goldwasser2021interactive, pmlr-v195-mutreja23a}. However, in informal domains, perfect completeness and soundness are unrealistic expectations for both models and humans. Therefore, we study relaxed versions of these properties: rather than perfect completeness, we expect the honest prover to achieve a high pass rate and the verifier to demonstrate high recall on honest solutions ($\varepsilon$-completeness); instead of perfect soundness, we expect the verifier to maintain a low false positive rate when evaluating adversarial inputs ($\varepsilon$-soundness).

\section{Simultaneous Play for Legibility}
In order to formalize the original PVG setup and justify our modifications to it, let us establish the baseline components and provide some definitions. \citet{kirchner2024prover} proposes to train (qualitatively) \textit{legible} provers and sound verifier using variants of the Prover-Verifier game \citep{anil2021learning}. They formulate the Prover-Verifier game as a simultaneous-move game between a role-conditioned prover $\pi_\theta$ and a verifier $v_\phi$. The prover, parameterized by $\theta$, generates solutions conditioned on its assigned role:

$$z_H = \pi_H(x, \text{helpful}), \quad z_S = \pi_S(x, \text{sneaky})$$

The verifier $v_\phi: \mathcal{X} \times \mathcal{Z} \to [-B, B]$ outputs a zero-centered convincingness score. The training proceeds via alternating optimization of two objectives. The verifier minimizes a classification loss:

$$L_V(\phi; \theta) = \frac{1}{2}\mathbb{E}_x\left[\ell(v_\phi(x, z_H), \corr(x, z_H))\right] + \frac{1}{2}\mathbb{E}_x\left[\ell(v_\phi(x, z_S), \corr(x, z_S))\right]$$

where $\ell$ is typically the cross-entropy loss or the Bradley-Terry loss. The prover maximizes a reward that encourages role-aligned behavior:

$$R_P(\theta; \phi) = \frac{1}{2}\mathbb{E}_x[r(v_\phi(x, z_H), \corr(x, z_H))] + \frac{1}{2}\mathbb{E}_x[r(v_\phi(x, z_S), 1 - \corr(x, z_S))]$$

The reward function $r(v, q)$ satisfies two critical properties:

\[\begin{aligned} \forall v_1, v_2 \in [0, 1] &: r(v_1, 1) > r(v_2, 0), \\ \forall v_1, v_2 \in [0, 1], v_1 > v_2 &: r(v_1, 1) > r(v_2, 1). \end{aligned}\]

\noindent where (i) states that role-aligned solutions always receive higher reward than misaligned ones, and (ii) states that among role-aligned solutions, higher verifier scores yield higher rewards.
\\
This formulation achieves qualitative empirical success: human evaluators rate helpful solutions as increasingly legible during training rounds, and the verifier demonstrates improved robustness against held-out adversarial examples. These properties suffice when the primary concern is unintentional errors, algebraic mistakes, logical non-sequiturs, or computational slips, but are not enough in our context.

\subsection{Limitations for Strategic Deception}
When facing intentionally deceptive agents in control scenarios, three pathologies emerge that fundamentally limit the applicability of simultaneous-play PVG:

\begin{itemize}
    \item \textbf{Style Leakage and Spurious Correlations.} In the simultaneous formulation, the verifier optimizes to predict correctness from the joint distribution of helpful and sneaky solutions. Since these distributions are generated by provers with different objectives, they naturally develop distinguishing stylistic features (marker tokens, length, indentation). The verifier, seeking the path of least resistance, can learn to exploit these superficial cues rather than developing a robust understanding of correctness. This is a weakness that \citet{kirchner2024prover} highlight as well in \S C of their paper.
    The simultaneous game can converge to a \textbf{weak equilibria}, resembling (implicit) \textbf{collusive dynamics} between the honest prover and the verifier. Note how the myopic objective of the players of the game does not allow us to ascribe intentionality to this behavior, but for lack of better term, we stick to this terminology.
    \item \textbf{Backdoor Blindness and Distribution Mismatch.} Most critically for control applications, the original formulation never exposes the verifier to the specific trigger inputs on which backdoored solutions misbehave. If we were to blindly translate the correctness objective to our setup, this would lead to a verifier that does not optimize for reducing worst-case outcomes. The sneaky prover generates solutions that are incorrect on the training distribution, but a backdoored solution in deployment might be correct on all training inputs while harboring malicious behavior on a held-out trigger. Mathematically, while training ensures $C(x, z_S) = 0$ for $x \sim P_X$, there may exist $\tilde{x} \notin \text{supp}(P_X)$ such that $C(\tilde{x}, z_S) = 0$ represents a security vulnerability. The training objective does not provide an incentive for the verifier to detect such backdoors, as they manifest themselves as correct solutions during training.\footnote{In \citet{kirchner2024prover}, sneaking in the training distribution is equivalent to sneaking in the test distribution, since being sneaky amounts to inserting a deliberate error anywhere in the reasoning trace. Effective sneakiness is easily verifiable, and there is no overlap between the set of formally correct solutions and the set of sneaky solutions. In code examples, this is not the case.}
\end{itemize}

\section{Sequential Play with Distribution Discrimination}

To address these limitations, we introduce three key modifications that transform the game from a legibility-focused framework to one suitable for adversarial control scenarios.

\paragraph{Sequential Stackelberg Structure.} We replace simultaneous play with a three-stage sequential game:
\begin{enumerate}
    \item The verifier $v_\phi$ commits to its evaluation strategy;
    \item The helpful prover $\pi_H(x)$ generates solution $z_H = \pi_H(x)$;
    \item The sneaky prover $\pi_S(x)$ observes $z_H$ and generates $z_S = \pi_S(x, z_H)$
\end{enumerate}

This sequential structure changes the game dynamics. The sneaky prover, observing the helpful solution, can copy any stylistic signals but must still satisfy its objective of introducing errors or backdoors.

\paragraph{Distribution Discrimination Objective.} Rather than training the verifier to predict correctness, we train it to identify which role distribution generated a solution. The verifier now aims to distinguish $P_{\text{honest}}$ from $P_{\text{sneaky}}$ rather than correct from incorrect solutions. Let
\[
\mathcal D_H=\bigl\{(x,z):z=\pi_H(x)\bigr\},\qquad
\mathcal D_S=\bigl\{(x,z):z=\pi_S(x)\bigr\}.
\]

The Kirchner verifier minimises the \textit{classification risk}

\[
R^\text{leg}(v)=\Pr_{(x,z)\sim\mathcal D_H\cup\mathcal D_S}
               \Bigl[v(x,z)\neq \corr(x,z)\Bigr],
\]

whereas our verifier minimises

\[
R^\text{ctl}(v)=\Pr_{(x,z)\sim\mathcal D_H\cup\mathcal D_S}
                \Bigl[v(x,z)\neq\mathbf 1[(x,z)\in\mathcal D_H]\Bigr].
\]

\noindent~While $\mathcal D_H$ may contain some incorrect samples (when $h$ errs) and $\mathcal D_S$ contains some correct samples (when the attack passes all unit test instances), the two risks may disagree on a non-empty set of points.  The solutions of the two ERM problems therefore differ unless $h$ is perfect, which is, in the limit, the outcome that $h$ aims to maximize.

\paragraph{Strength of distribution discrimination objective}
From the verifier's standpoint, we achieve $\varepsilon$-soundness via a modified Bradley-Terry objective:
$$L_V(\phi) = -\mathbb{E}_{x,z_H,z_S}\left[\log \sigma(v_\phi(x, z_H) - v_\phi(x, z_S))\right] + \lambda ||v_\phi||^2$$
\noindent~Whereas from the provers' point of view, the objective become:
$$R_H = v_\phi(x, z_H) \cdot \corr(x, z_H)$$
\\
while the sneaky reward incorporates backdoor success:
$$R_S = v_\phi(x, z_S) \cdot \success(x, z_S)$$\noindent
\smallskip
where we define the \emph{back-door success indicator}
\[
  \success(x,z_S):=
    \mathbf 1\bigl[\text{``$z_S$ embeds a trigger that changes behaviour
                     on some input''}\bigr].
\]
The sneaky prover’s utility will depend on this flag.\\
By learning to recognize the generative process rather than just the output property, the verifier is exposed to even stronger adversarial pressure by the sneaky prover, which in turn helps it develop more robust features that transfer across different instantiations of sneaky behavior.

\section{Theoretical Guarantees under Sequential Play}
Despite these substantial modifications, we can establish that the sequential game preserves the key theoretical properties of the original formulation.
\begin{definition}[Sequential Controllability]
A triple $(V, \dH, \dS)$ is $(\alpha, \beta)$-Sequentially-Controllable if the system is reliable and robust:
\begin{enumerate}
    \item (\textbf{$\alpha$-Completeness}) Both hold:
    \begin{align*}
        &\mathbb{E}_x\left[\corr(x, \pi_H(x))\right] \geq 1 - \alpha, \\
        &\mathbb{P}\left[v(x, \pi_H(x)) > v(x, \pi_S(x, \pi_H(x)))\right] \geq 1 - \alpha
    \end{align*}
    \item (\textbf{$\beta$-Soundness}) For all measurable $\pi_s': \mathcal{X} \times \mathcal{Z} \to \mathcal{Z}$,
    \[
        \mathbb{P}\left[
            \success(x, \pi_s'(x, \pi_H(x))) = 1
            \wedge\
            v(x, \pi_H(x)) \leq v(x, \pi_s'(x, \pi_H(x)))
        \right] \leq \beta
    \]
\end{enumerate}
\end{definition}

\begin{theorem}[Equivalence]
Under standard regularity conditions, an $(\varepsilon_V, \varepsilon_H, \varepsilon_S)$-approximate Stackelberg equilibrium of the sequential game is equivalent to an $(\alpha, \beta)$-sequentially controllable configuration with explicit moduli relating the parameters.
\begin{enumerate}
\item[(i)] Every $(\varepsilon_V,\varepsilon_H,\varepsilon_S)$-Stackelberg
           equilibrium is $(\alpha,\beta)$-sequentially controllable with
           \[
             \alpha=\max\!\Bigl(\tfrac{\varepsilon_H}{\gamma},
                                \tfrac{\varepsilon_V}{\ln 2}\Bigr),
             \qquad
             \beta=\tfrac{\varepsilon_V}{\ln 2}+\tfrac{\varepsilon_S}{\delta}.
           \]
\item[(ii)] Conversely, if $(V,\pi_H,\pi_S)$ is
           $(\alpha,\beta)$-controllable and $\pi_S$ is an
           $\varepsilon_S'$-best response, then it is an
           $(\varepsilon_V,\varepsilon_H,\varepsilon_S)$-equilibrium with
           \[
             \varepsilon_V=\ln 2+\alpha(C_B-\ln 2),\;
             \varepsilon_H=\max(\gamma,2BL)\alpha,\;
             \varepsilon_S=\Delta\beta.
           \]
\end{enumerate}
\end{theorem}\noindent
The proof, detailed in Appendix A, establishes that the sequential structure does not compromise the convergence guarantees while providing stronger robustness properties. The key insight is that the verifier's ability to distinguish distributions translates directly to detecting deceptive behavior, even from more capable future adversaries. Moreover, in Appendix B, we detail how the shift from simultaneous to sequential game does not affect the \emph{kind} of equilibria achievable by the game, but rather acts as a \emph{equilibrium refinement} mechanism.

\section{Addressing the Original Pathologies}
The sequential reformulation attempts to address the limitations of the original framework:

\paragraph{Style Leakage Mitigation.} Sequentiality forces stylistic alignment between honest and sneaky outputs. This eliminates the simple strategy of developing distinguishing stylistic markers, or, without ascribing intentionality, any strategy relying on spurious stylistic correlations. The verifier must learn features genuinely indicative of the generative process rather than superficial correlates.

\paragraph{Avoiding Weak Equilibria.} The sequential best-response structure ensures the honest prover cannot simply coordinate with the verifier on a shared signaling strategy. After observing $z_H$, the sneaky prover could copy any coordinated signal, but must still introduce actual errors or backdoors to achieve its reward. This forces the verifier to develop genuinely discriminative capabilities.

\paragraph{Backdoor Detection.} By explicitly incorporating trigger demonstrations in the training process, the verifier learns to recognize the patterns of backdoored solutions. The perturbation testing ensures robustness—a backdoor that only works on precisely $\tilde{x}$ but fails on slight variants receives low reward, encouraging the sneaky prover to create more robust backdoors that the verifier must learn to detect.
\\
This reformulation transforms the Prover-Verifier game from an elegant framework for improving legibility into a practical tool for adversarial robustness in control scenarios. The theoretical guarantees ensure convergence to meaningful equilibria, while the structural modifications address the specific challenges of defending against strategic deception.

\section{Research Questions}
With this in mind, we therefore pose 2 research questions:

\begin{tcolorbox}[
    colback=gray!10,
    colframe=black,
    boxrule=0.5pt,
    arc=3mm,
    left=10pt,
    right=10pt,
    top=8pt,
    bottom=8pt
]
\centering
\textit{\textbf{\normalsize Research Question \#1:}} \\[0.3em]
\textit{``Does a naïve implementation of the Prover-Verifier game lead to a good approximate equilibrium? That is, does the equilibrium that the players in the game converge to, or approximate convergence to, lead to the intended outcome?''}
\end{tcolorbox}\noindent
In this game, the incentives for the honest prover and the verifier are aligned; solution verification is a harder task than coordinating on a signaling cue for the verifier, and as long as the honest prover is using a consistent signal to flag its role in its rollouts (even inadvertently, we do not ascribe intentionality to myopically optimized systems), it is easier for the verifier to make its classification based on it rather than based on its correctness. We claim that this warrants our reformulation of the Prover-Verifier Game, requiring the new variant we present in this chapter.

\begin{tcolorbox}[
    colback=gray!10,
    colframe=black,
    boxrule=0.5pt,
    arc=3mm,
    left=10pt,
    right=10pt,
    top=8pt,
    bottom=8pt
]
\centering
\textit{\textbf{\normalsize Research Question \#2:}} \\[0.3em]
\textit{``Does simulating an attack policy via Prover-Verifier Games with trusted models make monitors resilient against more capable untrusted models?''}
\end{tcolorbox}\noindent

Say we are in an AI control scenario where we must instantiate a monitor $M$ for a protocol, and must choose between a weak trusted model $T$ and a more capable untrusted model $U$. Let us also assume for the sake of this protocol that the capabilities (or affordances) that these models possess have the following power relation:

$$\text{capabilities}_U \succ \text{capabilities}_T$$

This is an example of the type of problem AI control tries to solve, and a basic strategy in the literature \citep{greenblatt2023ai} looks like this: every time $U$ is used and it takes some action, use $M$ to check the action is OK before allowing execution. If the action is not OK, it's not executed (by $U$) and humans are alerted about suspicious activity.
\\
Currently, control protocols assume the monitor to leverage a mix of Human auditing $H$ (constrained to be rare, due to its cost), trusted monitoring (i.e., leveraging the trusted model $T$ as the monitor to provide an assessment of the risk of outputs coming from untrusted sources) and untrusted monitoring (i.e., for critical problems, leveraging untrusted policies to monitor untrusted outputs. Under some constraints and assumptions (statelessness, communication bottleneck, …), untrusted monitoring is often considered to be among the most promising protocols in terms of safety (i.e., probability of detecting worst-case outcomes) and usefulness (i.e., amount of useful labor that can be extracted), though we claim the protocol may be overly optimistic. Following \citet{greenblatt2023ai}, in order to leverage an untrusted model as a monitor, we would need to calibrate and test its ability to detect malicious outputs coming from the attacker/red team, an instance of itself. One way of probing for this involves honeypotting an untrusted model to see whether it would catch handcrafted misalignment. If this is the case, then we could be fairly confident in its judgement for actual outputs coming from $U$. A strong assumption made here is based on the fact that we can mimic almost perfectly the distribution of the attackers' outputs, in such a way that it would not be possible for the monitor to discriminate on the basis of the output being from a different distribution, that is, synthetic attacks may be so off-policy that they become easy to spot. What if the AI was \emph{situationally aware} in such a way that its verdict was on the basis of the synthetic inputs being blatantly synthetic rather than on their correctness?\footnote{Albeit speculative, situational awareness has been demonstrated to be increasingly stronger in frontier AI systems \citep{meinke2024frontier}.}
One way we claim can help us make better use of trusted monitors is by training them \emph{in simulation}.\\
Let us assume we have unlimited and unconstrained access only to $T$, and we know we ought to be prepared to evaluate a more capable model $U$ some time in the future. We can simulate $U$ by training attack policies via the aforementioned Prover-Verifier Game with the triplet of models $M, T, T_U$, where $M$ and $T$ stem from above and $T_U$ is a copy of the trusted model $T$ role-conditioned to be untrusted. The adversarial nature of the PVG allows us to train $T_U$ to mimic what a purposefully misaligned model $U$ would do, as well as gathering data for making $M$ more robust at detecting injections. The byproduct of this training procedure we care about is the monitor, for which we are going to assess its generalization capability in deployment against a stronger model explicitly tasked to backdoor.
